% !PRINTED 219
\weberpage{219}

\begin{center}
\S.\,71.

\gesperrt{Transformation der cubischen Gleichung.}
\end{center}

Der \gesperrt{Hermite'sche Satz} (\S.\,69) reicht für die cubische Gleichung aus, die Transformation ohne weitere Rechnung auszuführen.

Wir haben dazu, wenn wir homogene Variable $z_1, z_2$ anwenden, simultane Invarianten der cubischen Form

\weq{1}{f(z_1,z_2) = a_0z_1^3 + a_1z_1^2z_2 + a_2z_1z_2^2 + a_3z_2^3}

und der Linearform

\weq{2}{T = -t_0z_1 + t_1z_2}

zu bilden. Wenn wir aber die lineare Substitution

\weq{3}{\begin{gathered} z_1 = \alpha z_1' + \beta z_2' \\ z_2 = \gamma z_1' + \delta z_2' \end{gathered}}

auf $T$ anwenden, so ergiebt sich

\[ -t_0z_1 + t_1z_2 = -t_0'z_1' + t_1'z_2', \]

wenn

\[ \begin{gathered} t_0' = +t_0\alpha - t_1\gamma, \\ t_1' = -t_0\beta + t_1\delta, \end{gathered} \]

oder

\weq{4}{\begin{gathered} rt_1 = \alpha t_1' + \beta t_0' \\ rt_0 = \gamma t_1' + \delta t_0'. \end{gathered}}

Diese Substitution geht aber aus (3) hervor, wenn man $z_1, z_2, z_1', z_2'$ durch $rt_1, rt_0, t_1', t_0'$ ersetzt. Wenn nun $I(a,t_1,t_0)$ eine simultane Invariante von $f$ und $T$ ist, homogen und vom $\nu^{\mathrm{ten}}$ Grade in $t_0, t_1$ und vom Gewicht $\lambda$, also

\[ I(a',t_1',t_0') = r^\lambda I(a,t_1,t_0), \]

so folgt durch diese Vertauschung

\[ I(a',z_1',z_2') = r^{\lambda-\nu}I(a,z_1,z_2), \]

d.\,h. es ist $I(a,z_1,z_2)$ eine Covariante von $f$.

\gesperrt{Man erhält also alle simultanen Invarianten von $f$ und $T$ aus den Covarianten von $f$, wenn man darin $z_1, z_2$ durch $t_1, t_0$ ersetzt.}

Die Covarianten von $f$ haben wir aber in \S.\,62, 63 vollständig kennen gelernt.

Danach ist es leicht, die cubische Gleichung zu bilden, die sich aus der Substitution \S.\,68, (12)
%CONT>>
